\chapter{tekortschieting van documentatie}
\label{ch:tekortschieting_documentatie}

De complexiteit van EMAV-web, zoals besproken in de voorgaande secties, stelt hoge eisen aan de documentatie van het systeem. In de praktijk blijkt echter dat de traditionele methoden voor software-documentatie tekortschieten wanneer ze worden toegepast op een hybride omgeving waar wetenschappelijk onderzoek en software-engineering elkaar ontmoeten. Dit gebrek aan effectieve documentatie is een van de hoofdoorzaken van de kloof tussen de IT-afdeling en de onderzoekers.

\section{Het probleem van \textquotedbl Documentation Drift\textquotedbl}
Een van de meest hardnekkige problemen binnen het EMAV-project is de zogenaamde documentation drift: het fenomeen waarbij de geschreven documentatie en de feitelijke broncode na verloop van tijd uit elkaar groeien. Dit probleem is niet uniek voor EMAV. \textcite{tan2024outdatedDocumentation} tonen aan dat softwaredocumentatie verouderd kan raken wanneer verwijzingen naar code-elementen in de documentatie blijven bestaan nadat die elementen in de broncode gewijzigd of verwijderd zijn. Omdat de emissiemodellen regelmatig worden bijgewerkt op basis van nieuwe wetenschappelijke inzichten, moet de documentatie telkens handmatig worden gesynchroniseerd. In een omgeving met hoge tijdsdruk wordt dit vaak verwaarloosd, waardoor onderzoekers beslissingen baseren op verouderde informatie.

Binnen de C\#-codebase uit zich dit vaak in commentaren die niet langer overeenkomen met de berekeningslogica in de handlers of rekenmodules. De risico's hiervan zijn ook empirisch onderzocht: \textcite{ibrahim2012commentBugs} beschrijven dat verouderde comments ontwikkelaars op het verkeerde spoor kunnen zetten en dat patronen in commentaarupdates voorspellende waarde hebben voor toekomstige bugs. Wanneer een onderzoeker vraagt "hoe wordt parameter X berekend?", moet een programmeur vaak de code induiken omdat de externe documentatie (zoals Word-bestanden of Wiki-pagina's) niet langer betrouwbaar is. Dit proces is inefficiënt en verhoogt de kans op foutieve interpretaties van de wetenschappelijke modellen.

\section{Technische focus versus Domeinfocus}
Bestaande automatische documentatietools voor .NET, zoals DocFX of Swagger, zijn primair ontworpen voor softwareontwikkelaars. DocFX positioneert zich als een generator voor .NET API-documentatie op basis van broncode en XML-commentaren \autocite{docfx}. Swagger/OpenAPI beschrijft op zijn beurt API's in termen van paths, operaties, parameters, request bodies en responses \autocite{swaggerPathsOperations}. Deze tools excelleren daardoor in het weergeven van de technische interface (API-endpoints, klassenhiërarchieën, methodesiganturen), maar falen in het uitleggen van de intentie en de wiskundige onderbouwing van de logica. Dit sluit aan bij onderzoek naar API-documentatie: \textcite{meng2018apiDocumentation} benadrukken dat API-documentatie vooral moet beantwoorden aan de informatiebehoeften van softwareontwikkelaars.

Voor een onderzoeker is de informatie dat een methode CalculateEmission(double factor) heet en een Task<double> retourneert, van weinig waarde. Wat de onderzoeker nodig heeft, is inzicht in welke wetenschappelijke bronnen aan de basis liggen van de gebruikte coëfficiënten en hoe de formules onderling samenhangen. Vanuit Domain-Driven Design wordt dit spanningsveld herkend als een taal- en modelprobleem: de software moet een gedeelde domeintaal ondersteunen die door ontwikkelaars en domeinexperts begrepen wordt \autocite{Khononov2022}. De huidige documentatievormen zijn te "IT-centrisch" en bieden geen abstractieniveau dat aansluit bij de leefwereld van de wetenschapper.

\section{Gebrek aan interactiviteit en executabiliteit}
Traditionele documentatie is statisch. Een onderzoeker kan een PDF-document wel lezen, maar hij kan er niet mee "spelen". In een project als EMAV is het essentieel dat onderzoekers de impact van parameterwijzigingen direct kunnen observeren. Statische documentatie dwingt de onderzoeker in een passieve rol: hij moet de logica in de tekst begrijpen, deze mentaal vertalen naar de code, en vervolgens wachten op de IT-afdeling om een aanpassing door te voeren.

Het ontbreken van een executabel aspect in de documentatie betekent dat er geen "levende" link is tussen de uitleg en de uitvoering. Dit gebrek aan interactiviteit versterkt het black-box effect van de applicatie. De onderzoekers hebben nood aan een vorm van documentatie die niet alleen beschrijft wat de code doet, maar die hen ook in staat stelt om de logica te valideren door middel van directe interactie, zonder dat zij daarvoor de volledige ontwikkelomgeving hoeven te beheersen.
